%=============================================================================
\section{Quick Start：科研工作者的日常流程}
%=============================================================================

如果你只想尽快跑起来，这一节就够了。后面的章节是详细解释。

\subsection{一次典型的工作流}

假设你已经在本地写好了代码，想在集群上跑。整个过程只需要 5 步：

\textbf{第 1 步：把代码同步到集群}

在你的\textbf{本地电脑}上运行：

\begin{lstlisting}[language=bash]
rsync -azP --exclude='.venv' --exclude='__pycache__' \
  ./my_project/ your_utln@login-prod.pax.tufts.edu:~/my_project/
\end{lstlisting}

如果你已经配置了 SSH 别名（见连接章节），可以简写为：

\begin{lstlisting}[language=bash]
rsync -azP --exclude='.venv' --exclude='__pycache__' \
  ./my_project/ hpc:~/my_project/
\end{lstlisting}

\texttt{rsync} 只传有变化的文件，所以第二次开始会很快。

\textbf{第 2 步：SSH 登录集群}

\begin{lstlisting}[language=bash]
ssh hpc    # 或 ssh your_utln@login-prod.pax.tufts.edu
\end{lstlisting}

\textbf{第 3 步：提交作业}

假设你的项目里已经有一个 SLURM 脚本 \texttt{run.sh}：

\begin{lstlisting}[language=bash]
cd ~/my_project
sbatch run.sh
# 输出: Submitted batch job 347502
\end{lstlisting}

不知道怎么写 SLURM 脚本？这里是一个最小的模板，保存为 \texttt{run.sh}：

\begin{lstlisting}[language=slurm]
#!/bin/bash
#SBATCH -J my_job               # 作业名
#SBATCH --time=00-06:00:00      # 最多跑 6 小时
#SBATCH -p gpu                  # 使用 GPU 分区
#SBATCH --gres=gpu:1            # 1 块 GPU
#SBATCH --cpus-per-task=4       # 4 个 CPU 核
#SBATCH --mem=16g               # 16GB 内存
#SBATCH --output=%x.%j.out      # 标准输出日志
#SBATCH --error=%x.%j.err       # 错误日志

module load miniforge/25.3.0    # 加载 Conda
module load cuda/12.9.0         # 加载 CUDA
conda activate my_env           # 激活你的环境

python train.py                 # 运行你的程序
\end{lstlisting}

如果不需要 GPU，把 \texttt{-p gpu} 改成 \texttt{-p batch}，删掉 \texttt{--gres} 那行。

\textbf{第 4 步：查看作业状态}

\begin{lstlisting}[language=bash]
squeue --me                    # 看作业是在跑(R)还是排队(PD)
cat my_job.347502.out          # 看输出日志（347502 换成你的 Job ID）
\end{lstlisting}

提交之后你可以断开 SSH 去做别的事，作业会在后台继续跑。

\textbf{第 5 步：下载结果}

作业跑完后，在你的\textbf{本地电脑}上运行：

\begin{lstlisting}[language=bash]
rsync -azP hpc:~/my_project/results/ ./results/
\end{lstlisting}

\subsection{每天重复的部分}

上面 5 步里，第 2-4 步是你每天都会做的。最常见的循环是：

\begin{enumerate}[nosep]
  \item 本地改代码
  \item \texttt{rsync} 同步到集群（10 秒）
  \item \texttt{ssh hpc}，\texttt{sbatch run.sh}（5 秒）
  \item 过一会儿 \texttt{squeue --me} 看状态
  \item 跑完后 \texttt{rsync} 拉结果回来
\end{enumerate}

\subsection{遇到问题时怎么办}

\begin{table}[h]
\centering
\small
\begin{tabular}{lp{8cm}}
\toprule
情况 & 怎么做 \\
\midrule
作业一直排队（PD） & \texttt{sinfo} 看分区是否有空闲节点；减少资源请求；尝试加 \texttt{preempt} 分区 \\
作业运行失败 & 看 \texttt{.err} 日志文件找报错信息 \\
不确定环境对不对 & 开一个交互式节点调试：\texttt{srun -p batch -t 0-1:00:00 --mem=4g --pty bash} \\
配额满了 & \texttt{watch -n 5 quota -s} 监控；\texttt{du -sh \textasciitilde/.cache/*} 找大文件；\texttt{rm -rf \textasciitilde/.cache/uv} 清缓存 \\
想取消所有作业 & \texttt{scancel -u \$USER} \\
\bottomrule
\end{tabular}
\end{table}
